\chapter{Discussion}
\label{chap:discussion}

\section{What the Benchmark Actually Shows}

The benchmark shows that, despite a significant domain shift between hybrid simulation data and paired juxtacellular and intracellular ground-truth data, fpos can be predicted with reasonable error margins and variance explainability of the target variable. Conversely, fmiss cannot yet be predicted with the same level of robustness. 

However, this thesis does not claim to have identified a model that universally solves the issue of spike-sorting quality prediction. 

\section{Returning to the Hypotheses}

\textbf{H1 is supported.} Hybrid simulation data and paired juxtacellular and intracellular ground-truth data are significantly shifted, as shown in \cref{tab:domain-separability,fig:pca-domain-family}. The domain classifier, a simple logistic regression \parencite{10.1111/j.2517-6161.1958.tb00292.x}, achieves 0.9765 AUC with unit-level features and 0.9997 AUC when recording context is included, as reported in \cref{tab:domain-separability}. These results are also reflected in the benchmark, showing that transfer learning from hybrid simulation data to paired ground-truth data remains stable only if hybrid simulation units are filtered through a structure that reflects that of paired ground-truth data, as shown in \cref{tab:fpos-core-benchmark,tab:fpos-robustness,fig:fpos-stability}.

\textbf{H2 is supported for \fpos{}, but not supported for \fmiss{}.} The strongest \fpos{} stack, Waveform-augmented stack, leverages context successfully, as shown in \cref{tab:fpos-family-preference,tab:fpos-lofo-summary,fig:shap}. Additionally, the grouped feature importance in \cref{fig:shap} shows that the recording-context feature family has the highest average SHAP contribution \parencite{lundberg2017unifiedapproachinterpretingmodel}. On the other hand, models including context could not robustly predict \fmiss{}, as shown in \cref{tab:fmiss-summary,tab:fmiss-robustness,fig:fmiss-benchmark-stability,tab:fmiss-lofo-summary,fig:fmiss-family-heatmap}. Therefore, it cannot be supported that context about the recordings adds a meaningful predictive signal for the \fmiss{} target.

\textbf{H3 is supported.} \fpos{} and \fmiss{} are structurally different prediction targets. The difference in performance with the same feature package and model for both the paired-only raw model and paired-only context, shown in \cref{tab:fpos-core-benchmark,tab:fmiss-summary}, and the broader cross-target asymmetry shown in \cref{tab:target-protocol-summary,fig:target-asymmetry}, highlights it. Additionally, \fpos{} clearly accounts for a significant percentage of the target's variance, and its error margins are meaningful, as shown in \cref{tab:fpos-core-benchmark,tab:fpos-lofo-summary,fig:fpos-stability,fig:target-asymmetry}. On the other hand, \fmiss{} models, at best, break even in terms of variance, and their error margins are significantly larger, as shown in \cref{tab:fmiss-summary,tab:fmiss-lofo-summary,fig:fmiss-benchmark-stability,fig:target-asymmetry}.

\textbf{H4 is supported.} There are no universally better-performing models. In reality, for \fpos{}, the strongest model clearly depends on the dataset family and evaluation protocol. On the one hand, under the recording-disjoint protocol, the Trust-filtered unlabeled stack achieves the highest performance in \mae{}, seed wins, and mean rank, as shown in \cref{tab:fpos-robustness,fig:fpos-stability}. On the other hand, under the leave-one-family-out protocol, the Waveform-augmented stack is the best-performing one, being the best model for three of the five paired families, as shown in \cref{tab:fpos-family-preference,tab:fpos-lofo-summary,fig:fpos-family-heatmap}.

\section{Why the Targets Diverge}

The asymmetry between \fpos{} and \fmiss{} is not just an optimisation issue. \fpos{} is more directly exposed by the current feature package, which captures contamination-like failure through waveform, autocorrelogram, amplitude, and recording-context summaries. By contrast, \fmiss{} is more entangled with overlap, template mismatch, bursting, thresholding, and local dynamics that are only partially visible in the current features. This helps explain why \fpos{} reaches a meaningful predictive regime, whereas \fmiss{} remains much harder to model.

\section{Why Families Differ}

Family-level differences should be interpreted cautiously. The leave-one-family-out results show that model performance changes across paired families rather than remaining fixed, which means that a configuration that works well for one family does not necessarily remain the strongest for another. This is sufficient to show that cross-family generalisation is not uniform, even if the exact biological reasons likely differ from one family to another.

\subsection{BOYDEN and Bursting Neurons}

At its essence, BOYDEN's family \parencite{doi:10.1152/jn.00650.2017} was meant to represent bursting neurons and see how spike sorters perform on them. But these changes in firing pattern affect all the model's assumptions about amplitude, contamination, and spike-firing statistics, leading to their breakdown. Additionally, as seen in \cref{fig:shap}, the waveform-augmented model relies heavily on contamination features such as the autocorrelogram bin around 20 and \texttt{isi\_violation\_rate}, which is very affected, as shown in \cref{fig:family-profiles}, where the BOYDEN family's ACG peaks at bin 20 instead of falling to a trough, highlighting the high level of ISI violation. In summary, the models cannot fully rely on the assumptions they made about key features due to the biological nature of the recording family.

\section{Scientific Relevance}

This thesis provides valuable scientific insights to the spike-sorting field. It demonstrated that a significant portion of the target variance and meaningful error margins can be derived from \fpos{}, which is one of the core variables of the ground-truth accuracy metric, as shown in \cref{tab:fpos-core-benchmark,tab:fpos-lofo-summary,fig:fpos-stability}. This paves the way for further exploration of the metric and its potential application in post-sorting quality assessment. Additionally, one could argue that, if a sufficiently robust model were developed across diverse electrophysiological settings, brain regions, and neuron types, it could greatly improve how spike sorters detect and manage false positives. Such a model could enable sorters to estimate their error in real time after each unit is identified and to iteratively optimise for that metric. This thesis does not claim that its models are yet capable of doing so, but it was the first, to the best of our knowledge, to attempt to predict the \fpos{} target and extract meaningful information.

Furthermore, the incremental performance gains observed across multiple families through transfer learning from hybrid simulation data, supported by a trust-based filtering mechanism, might establish a precedent for future supervised learning methods in the spike-sorting field, as shown in \cref{tab:fpos-core-benchmark,tab:fpos-robustness,fig:fpos-stability,fig:fpos-family-heatmap}.

Finally, the structural difference between the \fmiss{} and \fpos{} targets indicates that false positives and missed spikes should not be seen as interchangeable parts of spike-sorting errors, as shown in \cref{tab:target-protocol-summary,fig:target-asymmetry}. This asymmetry implies that each target may need different feature datasets and modelling techniques. More broadly, the poor results obtained for \fmiss{} imply that current quality metrics may be more effective at identifying contamination than at detecting missing spikes, as shown in \cref{tab:fmiss-summary,tab:fmiss-lofo-summary,fig:fmiss-benchmark-stability}.

\section{Future work}

As highlighted in section 3.6, increasing the amount of paired ground truth data can lead to a significant improvement in performance. Due to the small size of the paired-labelled dataset, the most direct expected performance gain would be to increase its size. 

Additionally, adding new paired families and recordings would help the model generalise across a wider range of experimental settings. This would first increase the biological coverage of the dataset, introducing new brain regions, neuron types, and edge cases not yet represented. On top of that, expanding the range of experimental settings covered, such as recordings acquired with new probe types and sorters not yet included, would help the model generalise to a wider range of experimental settings.

Another ambition of this project would be to integrate the model into real post-sorting quality control pipelines. 

Finally, now that we are aware of the target asymmetry between fpos and fmiss, it would be relevant to see whether fmiss can be predicted using more ground-truth paired data and features centred on capturing overlap, bursting, and local dynamics between neurons.

\section{Limitations, Biological Scope, and Implications}

One of the main limitations is that the paired labelled set contains only 207 sorted units, and although they represent different brain regions, hardware simulations, sorters, and edge cases, they do not cover the full spectrum of potential electrophysiological recordings. Family coverage is uneven across sorters and recordings, meaning some families are more represented than others. This benchmark is therefore scientifically useful, though not exhaustive, due to the limited amount of paired juxtacellular and intracellular ground-truth data available online.

Another limitation is representational. The feature dataset is built from summaries of putative units rather than their raw traces, and from approximations of their waveforms and autocorrelograms rather than the exact ones. This is a deliberate design choice, as our entire dataset spans multiple terabytes of recordings, making it infeasible within our capabilities to store and leverage every piece of information. Additionally, we assume that the scenario in which our models are applied is post-sorting quality prediction.

Despite these limitations, the benchmark still supports practical recommendations. \fpos{} and \fmiss{} should be treated as separate deployment targets. Recording-disjoint evaluation should remain the default whenever context features are involved. Context-aware tree models are the strongest practical baseline family in this low-label tabular setting. Waveform enrichment is a regime-dependent gain rather than a guaranteed universal improvement.

The main next step is therefore not simply more tuning. For \fmiss{}, the priority is richer representations of missed-spike mechanisms. For \fpos{}, a promising direction is not a search for one universal final configuration, but a better understanding of when waveform augmentation is worth the extra complexity and when simpler structure-aware transfer remains safer.
